\documentclass[11pt, a4paper]{article}
\usepackage[a4paper, top=2.5cm, bottom=2.5cm, left=2cm, right=2cm]{geometry}
\usepackage{fontspec}
\usepackage[english, bidi=basic, provide=*]{babel}
\babelprovide[import, onchar=ids fonts]{english}
\babelfont{rm}{TeX Gyre Termes}
\babelfont{sf}{TeX Gyre Heros}
\babelfont{tt}{TeX Gyre Cursor}
\usepackage{enumitem}
\usepackage{amsmath}

\title{\textbf{Chapter 3: Water and Life \\ 75 Multiple Choice Questions}}
\author{}
\date{}

\begin{document}

\maketitle

\begin{enumerate}

% ---------------------------------------
% 3.1 Polar covalent bonds & H-bonds (1-12)
% ---------------------------------------

\item \textbf{In a single molecule of water, two hydrogen atoms are bonded to a single oxygen atom by:}
\begin{enumerate}[label=\Alph*.]
    \item Hydrogen bonds
    \item Nonpolar covalent bonds
    \item Polar covalent bonds
    \item Ionic bonds
    \item Van der Waals interactions
\end{enumerate}
\textbf{Answer:} C \\
\textbf{Explanation:} Oxygen is more electronegative than hydrogen, so it pulls shared electrons closer, creating polar covalent bonds within the molecule.

\item \textbf{The partial negative charge at one end of a water molecule is attracted to the partial positive charge of another water molecule. What is this attraction called?}
\begin{enumerate}[label=\Alph*.]
    \item A covalent bond
    \item A hydrogen bond
    \item An ionic bond
    \item A hydration shell
    \item A hydrophobic interaction
\end{enumerate}
\textbf{Answer:} B \\
\textbf{Explanation:} The attraction between oppositely charged regions of separate polar molecules is a hydrogen bond.

\item \textbf{Which element in a water molecule is the most electronegative?}
\begin{enumerate}[label=\Alph*.]
    \item Hydrogen
    \item Oxygen
    \item Carbon
    \item Nitrogen
    \item Helium
\end{enumerate}
\textbf{Answer:} B \\
\textbf{Explanation:} Oxygen is highly electronegative, causing unequal sharing of electrons in the water molecule.

\item \textbf{A single water molecule can form up to how many hydrogen bonds with neighboring water molecules?}
\begin{enumerate}[label=\Alph*.]
    \item 1
    \item 2
    \item 3
    \item 4
    \item 5
\end{enumerate}
\textbf{Answer:} D \\
\textbf{Explanation:} Each water molecule can hydrogen bond with four others (one at each hydrogen, two at the oxygen's lone electron pairs).

\item \textbf{The shape of a water molecule is most accurately described as:}
\begin{enumerate}[label=\Alph*.]
    \item Linear
    \item Tetrahedral
    \item A wide V
    \item Spherical
    \item Hexagonal
\end{enumerate}
\textbf{Answer:} C \\
\textbf{Explanation:} The bond angle between the two hydrogen atoms is 104.5 degrees, giving the molecule a wide V shape.

\item \textbf{In a liquid state, the hydrogen bonds between water molecules are:}
\begin{enumerate}[label=\Alph*.]
    \item Stronger than covalent bonds.
    \item Permanent and immobile.
    \item Fragile, constantly breaking and re-forming.
    \item Nonexistent.
    \item Repulsive.
\end{enumerate}
\textbf{Answer:} C \\
\textbf{Explanation:} Liquid water's dynamic nature is due to weak hydrogen bonds that form, break, and re-form very rapidly.

\item \textbf{Which type of bond must be broken for water to vaporize?}
\begin{enumerate}[label=\Alph*.]
    \item Ionic bonds
    \item Nonpolar covalent bonds
    \item Polar covalent bonds
    \item Hydrogen bonds
    \item Both polar covalent and hydrogen bonds
\end{enumerate}
\textbf{Answer:} D \\
\textbf{Explanation:} Changing from liquid to gas requires breaking the intermolecular hydrogen bonds, not the intramolecular covalent bonds.

\item \textbf{Water's polarity directly results from:}
\begin{enumerate}[label=\Alph*.]
    \item Its large molecular mass.
    \item The unequal sharing of electrons between oxygen and hydrogen.
    \item The equal sharing of electrons.
    \item The loss of electrons to the environment.
    \item Its inability to form covalent bonds.
\end{enumerate}
\textbf{Answer:} B \\
\textbf{Explanation:} Unequal sharing caused by oxygen's high electronegativity creates partial positive and negative poles.

\item \textbf{Which of the following is true about the strength of hydrogen bonds in water?}
\begin{enumerate}[label=\Alph*.]
    \item They are twice as strong as covalent bonds.
    \item They are equal in strength to covalent bonds.
    \item They are about 1/20th as strong as covalent bonds.
    \item They are stronger than ionic bonds in a dry crystal.
    \item They cannot be broken by heat.
\end{enumerate}
\textbf{Answer:} C \\
\textbf{Explanation:} Hydrogen bonds are relatively weak intermolecular forces compared to intramolecular covalent bonds.

\item \textbf{If water were a nonpolar molecule, how would its state at room temperature likely change?}
\begin{enumerate}[label=\Alph*.]
    \item It would be a solid.
    \item It would be a gas.
    \item It would remain a liquid but become highly viscous.
    \item It would become highly acidic.
    \item It would spontaneously combust.
\end{enumerate}
\textbf{Answer:} B \\
\textbf{Explanation:} Without hydrogen bonds to hold molecules together, water's low molecular mass would cause it to be a gas at room temperature, like methane.

\item \textbf{The partial charges on a water molecule are designated by which Greek letter?}
\begin{enumerate}[label=\Alph*.]
    \item Alpha ($\alpha$)
    \item Beta ($\beta$)
    \item Gamma ($\gamma$)
    \item Delta ($\delta$)
    \item Sigma ($\sigma$)
\end{enumerate}
\textbf{Answer:} D \\
\textbf{Explanation:} Delta minus ($\delta-$) represents partial negative and delta plus ($\delta+$) represents partial positive.

\item \textbf{In an aqueous solution, water molecules are attracted to:}
\begin{enumerate}[label=\Alph*.]
    \item Only other water molecules.
    \item Nonpolar molecules like oils.
    \item Both other polar molecules and ions.
    \item Only anions.
    \item Only cations.
\end{enumerate}
\textbf{Answer:} C \\
\textbf{Explanation:} Water's polarity allows it to form hydrogen bonds with other polar substances and electrostatic attractions with both positive and negative ions.

% ---------------------------------------
% 3.2 Property 1: Cohesion & Adhesion (13-20)
% ---------------------------------------

\item \textbf{The phenomenon where water molecules are held together by hydrogen bonds is known as:}
\begin{enumerate}[label=\Alph*.]
    \item Adhesion
    \item Cohesion
    \item Surface tension
    \item Specific heat
    \item Evaporation
\end{enumerate}
\textbf{Answer:} B \\
\textbf{Explanation:} Cohesion is the linking together of like molecules.

\item \textbf{Water clings to the cell walls of plant vessels. This is an example of:}
\begin{enumerate}[label=\Alph*.]
    \item Cohesion
    \item Adhesion
    \item High specific heat
    \item Heat of vaporization
    \item Hydration
\end{enumerate}
\textbf{Answer:} B \\
\textbf{Explanation:} Adhesion is the clinging of one substance to another, unlike substance (e.g., water to plant cell walls).

\item \textbf{Which property allows water to travel upwards against gravity in a plant?}
\begin{enumerate}[label=\Alph*.]
    \item Adhesion alone
    \item High specific heat
    \item Cohesion and adhesion
    \item Lower density as a solid
    \item High heat of vaporization
\end{enumerate}
\textbf{Answer:} C \\
\textbf{Explanation:} Cohesion pulls the water column up, while adhesion to the plant wall resists the downward pull of gravity.

\item \textbf{A water strider bug can walk across the surface of a pond because of water's high:}
\begin{enumerate}[label=\Alph*.]
    \item Surface tension
    \item Kinetic energy
    \item Specific heat
    \item Solvent capacity
    \item Density
\end{enumerate}
\textbf{Answer:} A \\
\textbf{Explanation:} High surface tension, a result of cohesive hydrogen bonding at the surface, resists external forces.

\item \textbf{Surface tension is closely related to which other property of water?}
\begin{enumerate}[label=\Alph*.]
    \item pH
    \item Cohesion
    \item Evaporative cooling
    \item Hydrophobic exclusion
    \item Molarity
\end{enumerate}
\textbf{Answer:} B \\
\textbf{Explanation:} Both surface tension and cohesion result directly from the hydrogen bonds linking water molecules together.

\item \textbf{If a chemical was added to water that disrupted hydrogen bonds, what would happen to its surface tension?}
\begin{enumerate}[label=\Alph*.]
    \item It would drastically increase.
    \item It would dramatically decrease.
    \item It would remain exactly the same.
    \item It would fluctuate randomly.
    \item The water would freeze.
\end{enumerate}
\textbf{Answer:} B \\
\textbf{Explanation:} Surface tension relies on hydrogen bonds; disrupting them would break the "film" on the water's surface.

\item \textbf{What is the primary physical force that cohesion and adhesion must overcome in a tall tree?}
\begin{enumerate}[label=\Alph*.]
    \item Wind
    \item Atmospheric pressure
    \item Gravity
    \item Sunlight
    \item Cellular respiration
\end{enumerate}
\textbf{Answer:} C \\
\textbf{Explanation:} To move water from roots to leaves, the water column must resist the downward pull of gravity.

\item \textbf{At the interface between water and air, the hydrogen bonds of water molecules are:}
\begin{enumerate}[label=\Alph*.]
    \item Absent.
    \item Ordered, producing high surface tension.
    \item Repelling the air molecules via ionic forces.
    \item Rapidly evaporating.
    \item Weaker than in the bulk liquid.
\end{enumerate}
\textbf{Answer:} B \\
\textbf{Explanation:} Surface molecules hydrogen-bond to those below and beside them, creating an invisible, ordered "film".

% ---------------------------------------
% 3.2 Property 2: Temperature & Heat (21-36)
% ---------------------------------------

\item \textbf{The energy of motion in atoms or molecules is defined as:}
\begin{enumerate}[label=\Alph*.]
    \item Potential energy
    \item Thermal energy
    \item Kinetic energy
    \item Specific heat
    \item Heat of vaporization
\end{enumerate}
\textbf{Answer:} C \\
\textbf{Explanation:} Kinetic energy is the energy associated with the movement of objects or particles.

\item \textbf{Temperature is fundamentally a measure of:}
\begin{enumerate}[label=\Alph*.]
    \item The total kinetic energy of a body of matter.
    \item The average kinetic energy of molecules in a body of matter.
    \item The amount of light absorbed by a liquid.
    \item The volume of a liquid.
    \item The total number of hydrogen bonds in a solution.
\end{enumerate}
\textbf{Answer:} B \\
\textbf{Explanation:} Temperature measures the average kinetic energy, while thermal energy measures the total kinetic energy.

\item \textbf{Which has more thermal energy: a hot cup of coffee or a cold swimming pool?}
\begin{enumerate}[label=\Alph*.]
    \item The coffee, because it has a higher temperature.
    \item The swimming pool, because it has a vastly larger volume.
    \item They have the same thermal energy.
    \item Thermal energy cannot be determined without measuring pH.
    \item The coffee, because it is more dense.
\end{enumerate}
\textbf{Answer:} B \\
\textbf{Explanation:} Though the coffee has higher temperature (average kinetic energy), the pool's massive volume means it contains far more total thermal energy.

\item \textbf{The amount of heat it takes to raise the temperature of 1 gram of water by 1 degree Celsius is one:}
\begin{enumerate}[label=\Alph*.]
    \item Joule
    \item Watt
    \item Calorie
    \item Kelvin
    \item Mole
\end{enumerate}
\textbf{Answer:} C \\
\textbf{Explanation:} A calorie (cal) is the standard unit used to define the specific heat of water.

\item \textbf{Water has a high specific heat because of its:}
\begin{enumerate}[label=\Alph*.]
    \item Low molecular mass
    \item Hydrogen bonds
    \item Covalent bonds
    \item Adhesion
    \item Ionic composition
\end{enumerate}
\textbf{Answer:} B \\
\textbf{Explanation:} Heat must first be absorbed to break hydrogen bonds before the water molecules can move faster to increase temperature.

\item \textbf{When water cools, what happens at the molecular level?}
\begin{enumerate}[label=\Alph*.]
    \item Covalent bonds are broken.
    \item Hydrogen bonds are broken, absorbing heat.
    \item Hydrogen bonds form, releasing heat.
    \item Molecules move much faster.
    \item The water splits into hydrogen and oxygen gas.
\end{enumerate}
\textbf{Answer:} C \\
\textbf{Explanation:} When water temperature drops, many hydrogen bonds form, which releases a significant amount of heat, slowing the cooling process.

\item \textbf{Which characteristic of water allows coastal areas to have milder climates than inland areas?}
\begin{enumerate}[label=\Alph*.]
    \item Low specific heat
    \item High specific heat
    \item Low heat of vaporization
    \item High surface tension
    \item Excellent solvent properties
\end{enumerate}
\textbf{Answer:} B \\
\textbf{Explanation:} Water's high specific heat allows oceans to absorb heat during the day/summer and release it at night/winter, buffering climate changes.

\item \textbf{The quantity of heat a liquid must absorb for 1 gram of it to be converted from liquid to the gaseous state is:}
\begin{enumerate}[label=\Alph*.]
    \item Specific heat
    \item Heat of fusion
    \item Heat of vaporization
    \item Thermal capacity
    \item Evaporative potential
\end{enumerate}
\textbf{Answer:} C \\
\textbf{Explanation:} Heat of vaporization is the energy required to vaporize a liquid; for water, it is unusually high.

\item \textbf{Evaporative cooling occurs because:}
\begin{enumerate}[label=\Alph*.]
    \item The molecules with the greatest kinetic energy leave as a gas.
    \item Water absorbs cold from the atmosphere.
    \item Hydrogen bonds release heat when they break.
    \item The remaining molecules suddenly accelerate.
    \item Sweating produces endothermic chemical reactions.
\end{enumerate}
\textbf{Answer:} A \\
\textbf{Explanation:} As the "hottest" (fastest) molecules evaporate, the average kinetic energy (temperature) of the remaining liquid decreases.

\item \textbf{Sweating helps humans cool down primarily due to water's:}
\begin{enumerate}[label=\Alph*.]
    \item Low molecular weight
    \item High surface tension
    \item High heat of vaporization
    \item Cohesion
    \item Neutral pH
\end{enumerate}
\textbf{Answer:} C \\
\textbf{Explanation:} Evaporation of sweat consumes a large amount of body heat because of water's high heat of vaporization.

\item \textbf{If water had a low specific heat, what would be the most likely consequence for organisms?}
\begin{enumerate}[label=\Alph*.]
    \item They would require less water to survive.
    \item They would easily boil or freeze with minor environmental temperature changes.
    \item Their blood would flow faster.
    \item They would not be able to sweat.
    \item They would float on water more easily.
\end{enumerate}
\textbf{Answer:} B \\
\textbf{Explanation:} Organisms are mostly water; a low specific heat would make them highly susceptible to lethal internal temperature fluctuations.

\item \textbf{Why does steam burns cause more severe damage than boiling water at the same temperature?}
\begin{enumerate}[label=\Alph*.]
    \item Steam contains acidic ions.
    \item Steam molecules move slower.
    \item Steam releases the massive heat of vaporization when it condenses on the skin.
    \item Steam is denser than liquid water.
    \item Steam exerts higher atmospheric pressure.
\end{enumerate}
\textbf{Answer:} C \\
\textbf{Explanation:} Condensing steam releases the massive amount of energy (heat of vaporization) it originally absorbed to become a gas.

\item \textbf{Compared to an equal mass of alcohol, water requires \_\_\_\_\_\_\_ heat to increase its temperature by 1 degree.}
\begin{enumerate}[label=\Alph*.]
    \item Less
    \item More
    \item The exact same amount of
    \item Zero
    \item Ten times less
\end{enumerate}
\textbf{Answer:} B \\
\textbf{Explanation:} Water has a much higher specific heat than most other substances, including alcohol.

\item \textbf{The SI unit of energy used alongside the calorie is the:}
\begin{enumerate}[label=\Alph*.]
    \item Watt
    \item Newton
    \item Joule
    \item Kelvin
    \item Dalton
\end{enumerate}
\textbf{Answer:} C \\
\textbf{Explanation:} One joule (J) equals 0.239 calories, and it is a standard scientific unit for energy.

\item \textbf{When a pot of water is heated on a stove, what must happen first before the temperature visibly rises?}
\begin{enumerate}[label=\Alph*.]
    \item Covalent bonds must break.
    \item Hydrogen bonds must break.
    \item The water must begin evaporating.
    \item The pot must melt.
    \item Surface tension must double.
\end{enumerate}
\textbf{Answer:} B \\
\textbf{Explanation:} Most of the initial heat energy is consumed in disrupting the hydrogen bonds before the molecules can move faster.

\item \textbf{Which process is responsible for preventing a tropical plant's leaves from overheating in the sun?}
\begin{enumerate}[label=\Alph*.]
    \item Surface tension
    \item Ice flotation
    \item Evaporative cooling (transpiration)
    \item Adhesion
    \item Hydration shelling
\end{enumerate}
\textbf{Answer:} C \\
\textbf{Explanation:} Evaporation of water from the leaves removes excess heat from the plant.

% ---------------------------------------
% 3.2 Property 3: Floating Ice (37-46)
% ---------------------------------------

\item \textbf{Water reaches its greatest density at what temperature?}
\begin{enumerate}[label=\Alph*.]
    \item $0^\circ\text{C}$
    \item $4^\circ\text{C}$
    \item $10^\circ\text{C}$
    \item $25^\circ\text{C}$
    \item $100^\circ\text{C}$
\end{enumerate}
\textbf{Answer:} B \\
\textbf{Explanation:} As water cools below 4 degrees Celsius, it begins to expand, reaching its minimum density as solid ice at 0 degrees Celsius.

\item \textbf{Why does ice float on liquid water?}
\begin{enumerate}[label=\Alph*.]
    \item Ice is strictly nonpolar.
    \item Ice repels water molecules.
    \item Hydrogen bonds stabilize and keep the molecules further apart in ice than in liquid water.
    \item The covalent bonds in ice expand.
    \item Ice absorbs air bubbles when freezing.
\end{enumerate}
\textbf{Answer:} C \\
\textbf{Explanation:} In ice, hydrogen bonds lock molecules into a crystalline lattice that is roughly 10\% less dense than the liquid phase.

\item \textbf{If ice were denser than liquid water, what would be the ecological consequence?}
\begin{enumerate}[label=\Alph*.]
    \item Oceans would not exist.
    \item All bodies of water would eventually freeze solid from the bottom up, killing most aquatic life.
    \item Aquatic animals would thrive in the cooler waters.
    \item The specific heat of water would drop to zero.
    \item Evaporation would stop entirely.
\end{enumerate}
\textbf{Answer:} B \\
\textbf{Explanation:} If ice sank, it would not insulate the water below, causing entire lakes and oceans to freeze solid over time.

\item \textbf{In ice, each water molecule forms how many stable hydrogen bonds with its neighbors?}
\begin{enumerate}[label=\Alph*.]
    \item 1
    \item 2
    \item 3
    \item 4
    \item 6
\end{enumerate}
\textbf{Answer:} D \\
\textbf{Explanation:} In the solid crystalline lattice, each water molecule is maximally hydrogen-bonded to exactly four other molecules.

\item \textbf{Which environment represents a habitat preserved due to the insulating property of floating ice?}
\begin{enumerate}[label=\Alph*.]
    \item Desert sand dunes
    \item Tropical rainforest canopy
    \item Liquid water below the frozen surface of a lake
    \item Geothermal hot springs
    \item Volcanic vents
\end{enumerate}
\textbf{Answer:} C \\
\textbf{Explanation:} Floating ice acts as a blanket, preventing the deep water below from freezing and allowing aquatic life to survive winter.

\item \textbf{As water warms from $0^\circ\text{C}$ to $4^\circ\text{C}$, it:}
\begin{enumerate}[label=\Alph*.]
    \item Expands and becomes less dense.
    \item Contracts and becomes more dense.
    \item Vaporizes immediately.
    \item Breaks all its covalent bonds.
    \item Forms a rigid crystalline lattice.
\end{enumerate}
\textbf{Answer:} B \\
\textbf{Explanation:} Between 0 and 4 degrees, the breaking of the rigid ice lattice allows molecules to pack closer together, increasing density.

\item \textbf{Global warming is severely affecting which animal that relies on intact Arctic sea ice for hunting?}
\begin{enumerate}[label=\Alph*.]
    \item Parrots
    \item Desert tortoises
    \item Ringed seals and polar bears
    \item Great white sharks
    \item Galápagos finches
\end{enumerate}
\textbf{Answer:} C \\
\textbf{Explanation:} Arctic sea ice provides a necessary platform for species like ringed seals and polar bears to hunt and survive.

\item \textbf{When comparing a gram of liquid water at $4^\circ\text{C}$ and a gram of ice at $0^\circ\text{C}$, the ice has:}
\begin{enumerate}[label=\Alph*.]
    \item More molecules
    \item Fewer molecules
    \item The same number of molecules but takes up more volume
    \item The same number of molecules but takes up less volume
    \item Higher kinetic energy
\end{enumerate}
\textbf{Answer:} C \\
\textbf{Explanation:} Density = mass/volume. Since mass is the same and ice is less dense, it must occupy more volume.

\item \textbf{During the spring turnover in a lake, the sinking of surface water is caused by:}
\begin{enumerate}[label=\Alph*.]
    \item Water reaching $4^\circ\text{C}$ and becoming densest.
    \item Water freezing at $0^\circ\text{C}$.
    \item Evaporation of surface water.
    \item Acid rain lowering the pH.
    \item Fish swimming to the bottom.
\end{enumerate}
\textbf{Answer:} A \\
\textbf{Explanation:} As surface ice melts and warms to 4 degrees, it becomes denser and sinks, mixing the lake's nutrients and oxygen.

\item \textbf{What structural feature of water prevents it from packing tightly when frozen?}
\begin{enumerate}[label=\Alph*.]
    \item The high mass of the oxygen atom.
    \item The fixed, rigid, "arm's-length" hydrogen bonds.
    \item The constant breaking of covalent bonds.
    \item The nonpolar nature of the molecule.
    \item The lack of electrons in hydrogen.
\end{enumerate}
\textbf{Answer:} B \\
\textbf{Explanation:} The hydrogen bonds lock the molecules into a spacious, 3D crystalline lattice that pushes molecules apart.

% ---------------------------------------
% 3.2 Property 4: Solvent of Life (47-58)
% ---------------------------------------

\item \textbf{A completely homogeneous mixture of two or more substances is called a:}
\begin{enumerate}[label=\Alph*.]
    \item Solute
    \item Solvent
    \item Solution
    \item Compound
    \item Suspension
\end{enumerate}
\textbf{Answer:} C \\
\textbf{Explanation:} A solution is a uniform mixture, like sugar fully dissolved in water.

\item \textbf{In a glass of salt water, the water is the \_\_\_\_\_\_\_\_ and the salt is the \_\_\_\_\_\_\_\_.}
\begin{enumerate}[label=\Alph*.]
    \item Solute; solvent
    \item Solvent; solute
    \item Solution; solute
    \item Solvent; solution
    \item Ion; compound
\end{enumerate}
\textbf{Answer:} B \\
\textbf{Explanation:} The dissolving agent is the solvent (water), and the substance being dissolved is the solute (salt).

\item \textbf{When an ionic compound like NaCl is dissolved in water, each ion is surrounded by a sphere of water molecules called a:}
\begin{enumerate}[label=\Alph*.]
    \item Crystalline lattice
    \item Hydration shell
    \item Cohesion circle
    \item Buffer zone
    \item Hydrophobic drop
\end{enumerate}
\textbf{Answer:} B \\
\textbf{Explanation:} Water's partial charges attract the ions, surrounding them in a hydration shell that separates and dissolves the salt.

\item \textbf{Which parts of water molecules point toward a dissolved sodium ($Na^+$) cation?}
\begin{enumerate}[label=\Alph*.]
    \item The partially positive hydrogen atoms
    \item The partially negative oxygen atoms
    \item Both hydrogen and oxygen atoms equally
    \item Neither, water repels cations
    \item The unbonded neutrons
\end{enumerate}
\textbf{Answer:} B \\
\textbf{Explanation:} The partially negative oxygen regions are electrostatically attracted to the positive sodium cation.

\item \textbf{A substance that has an affinity for water is termed:}
\begin{enumerate}[label=\Alph*.]
    \item Hydrophobic
    \item Nonpolar
    \item Hydrophilic
    \item Hydrated
    \item Acidic
\end{enumerate}
\textbf{Answer:} C \\
\textbf{Explanation:} Hydrophilic (water-loving) substances are usually polar or ionic and readily interact with or dissolve in water.

\item \textbf{Substances that are nonionic and nonpolar, and seem to repel water, are termed:}
\begin{enumerate}[label=\Alph*.]
    \item Hydrophilic
    \item Solvents
    \item Aqueous
    \item Hydrophobic
    \item Basic
\end{enumerate}
\textbf{Answer:} D \\
\textbf{Explanation:} Hydrophobic (water-fearing) substances, like oils, do not mix with water because they cannot form hydrogen bonds.

\item \textbf{Large protein molecules can dissolve in water if they have \_\_\_\_\_\_\_ regions on their surface.}
\begin{enumerate}[label=\Alph*.]
    \item Ionic and polar
    \item Strictly nonpolar
    \item Hydrophobic
    \item Uncharged
    \item Gaseous
\end{enumerate}
\textbf{Answer:} A \\
\textbf{Explanation:} Large molecules like proteins (e.g., lysozyme) can dissolve if they possess enough polar and ionic regions to form hydrogen bonds with water.

\item \textbf{Cellulose, found in plant cell walls, is hydrophilic but does not dissolve in water. Why?}
\begin{enumerate}[label=\Alph*.]
    \item It is made of nonpolar covalent bonds.
    \item It is hydrophobic.
    \item Its molecules are too large to dissolve, despite adhering to water.
    \item It is highly acidic and reacts with water.
    \item It turns water into a solid.
\end{enumerate}
\textbf{Answer:} C \\
\textbf{Explanation:} Cellulose contains partial charges that attract water (hence it absorbs water, like in a towel), but it is too massive to completely break apart and dissolve.

\item \textbf{The sum of the masses of all atoms in a single molecule is its:}
\begin{enumerate}[label=\Alph*.]
    \item Molarity
    \item Molecular mass
    \item Atomic number
    \item Specific heat
    \item Avogadro count
\end{enumerate}
\textbf{Answer:} B \\
\textbf{Explanation:} By multiplying the number of each type of atom by its atomic mass and adding them, you get the molecular mass.

\item \textbf{One mole (mol) of any substance contains exactly how many molecules?}
\begin{enumerate}[label=\Alph*.]
    \item $1.0 \times 10^{14}$
    \item $6.02 \times 10^{23}$
    \item $3.14 \times 10^{15}$
    \item $1.0 \times 10^7$
    \item $9.8 \times 10^{10}$
\end{enumerate}
\textbf{Answer:} B \\
\textbf{Explanation:} Avogadro's number ($6.02 \times 10^{23}$) defines a mole, allowing chemists to relate molecular mass to grams.

\item \textbf{How many grams of sucrose ($C_{12}H_{22}O_{11}$) are in 1 mole of sucrose, given its molecular mass is 342 daltons?}
\begin{enumerate}[label=\Alph*.]
    \item 1 gram
    \item $6.02 \times 10^{23}$ grams
    \item 342 grams
    \item 12 grams
    \item 100 grams
\end{enumerate}
\textbf{Answer:} C \\
\textbf{Explanation:} A mole of a substance has a mass in grams precisely equal to its molecular mass in daltons.

\item \textbf{Molarity (M) is defined as the number of:}
\begin{enumerate}[label=\Alph*.]
    \item Moles of solvent per liter of solute.
    \item Moles of solute per liter of solution.
    \item Grams of solute per liter of solution.
    \item Molecules of solute per milliliter.
    \item Moles of solute per kilogram of solvent.
\end{enumerate}
\textbf{Answer:} B \\
\textbf{Explanation:} Molarity is the standard unit of concentration used by biologists for aqueous solutions.

% ---------------------------------------
% 3.3 Acids, Bases, and pH (59-69)
% ---------------------------------------

\item \textbf{When a hydrogen ion ($H^+$) transfers from one water molecule to another, the water molecule that lost the proton becomes a(n):}
\begin{enumerate}[label=\Alph*.]
    \item Hydronium ion ($H_3O^+$)
    \item Hydroxide ion ($OH^-$)
    \item Hydrogen gas ($H_2$)
    \item Oxygen gas ($O_2$)
    \item Peroxide ion
\end{enumerate}
\textbf{Answer:} B \\
\textbf{Explanation:} Losing an $H^+$ leaves behind an oxygen and one hydrogen with an extra electron, making it $OH^-$.

\item \textbf{The water molecule that receives the extra proton ($H^+$) becomes a:}
\begin{enumerate}[label=\Alph*.]
    \item Hydroxide ion ($OH^-$)
    \item Hydrogen atom (H)
    \item Hydronium ion ($H_3O^+$)
    \item Hydration shell
    \item Buffer
\end{enumerate}
\textbf{Answer:} C \\
\textbf{Explanation:} The $H^+$ binds to the lone pair of another water molecule to form $H_3O^+$.

\item \textbf{In pure water at $25^\circ\text{C}$, the concentration of $H^+$ ions is:}
\begin{enumerate}[label=\Alph*.]
    \item $10^{-14}$ M
    \item $10^{-7}$ M
    \item $1$ M
    \item $7$ M
    \item $0$ M
\end{enumerate}
\textbf{Answer:} B \\
\textbf{Explanation:} Pure water is neutral, meaning $[H^+] = [OH^-] = 10^{-7}$ M.

\item \textbf{A substance that increases the hydrogen ion concentration of a solution is a(n):}
\begin{enumerate}[label=\Alph*.]
    \item Base
    \item Solvent
    \item Buffer
    \item Acid
    \item Isotope
\end{enumerate}
\textbf{Answer:} D \\
\textbf{Explanation:} Acids donate $H^+$ to the solution, increasing the overall $[H^+]$.

\item \textbf{A substance that reduces the hydrogen ion concentration of a solution is a(n):}
\begin{enumerate}[label=\Alph*.]
    \item Acid
    \item Base
    \item Solvent
    \item Solute
    \item Catalyst
\end{enumerate}
\textbf{Answer:} B \\
\textbf{Explanation:} Bases reduce $[H^+]$ either by directly accepting $H^+$ or by dissociating to form $OH^-$ which then reacts with $H^+$ to form water.

\item \textbf{Hydrochloric acid (HCl) is a strong acid because it:}
\begin{enumerate}[label=\Alph*.]
    \item Reversibly releases and accepts $H^+$.
    \item Completely dissociates into $H^+$ and $Cl^-$ in water.
    \item Acts as a buffer in human blood.
    \item Has a high pH.
    \item Combines with water to form a solid.
\end{enumerate}
\textbf{Answer:} B \\
\textbf{Explanation:} Strong acids completely dissociate in aqueous solutions, irreversibly contributing all their $H^+$ to the solution.

\item \textbf{Ammonia ($NH_3$) acts as a base when it dissolves in water because it:}
\begin{enumerate}[label=\Alph*.]
    \item Donates $OH^-$ directly to the water.
    \item Accepts an $H^+$ from water to become an ammonium ion ($NH_4^+$).
    \item Increases the $[H^+]$ concentration.
    \item Completely dissociates into nitrogen and hydrogen.
    \item Acts as a strong acid.
\end{enumerate}
\textbf{Answer:} B \\
\textbf{Explanation:} Ammonia binds with $H^+$ ions from the solution, thereby reducing the $H^+$ concentration and acting as a base.

\item \textbf{The ion product of water at $25^\circ\text{C}$ is mathematically expressed as:}
\begin{enumerate}[label=\Alph*.]
    \item $[H^+] + [OH^-] = 14$
    \item $[H^+] \times [OH^-] = 10^{-14}$
    \item $[H^+] / [OH^-] = 1$
    \item $\log[H^+] = 7$
    \item $[H^+] = 10^{-14}$
\end{enumerate}
\textbf{Answer:} B \\
\textbf{Explanation:} The product of the $H^+$ and $OH^-$ concentrations is constant at $10^{-14}$ in any aqueous solution.

\item \textbf{The pH of a solution is defined by the negative logarithm of the:}
\begin{enumerate}[label=\Alph*.]
    \item Hydroxide ion concentration
    \item Hydrogen ion concentration
    \item Total water volume
    \item Buffer capacity
    \item Molar mass
\end{enumerate}
\textbf{Answer:} B \\
\textbf{Explanation:} pH = $-\log[H^+]$.

\item \textbf{If a solution has a pH of 3, it is:}
\begin{enumerate}[label=\Alph*.]
    \item Neutral
    \item Strongly basic
    \item Acidic
    \item A buffer
    \item Pure water
\end{enumerate}
\textbf{Answer:} C \\
\textbf{Explanation:} A pH value less than 7 indicates an acidic solution.

\item \textbf{A solution with a pH of 4 has how many times more hydrogen ions ($H^+$) than a solution with a pH of 6?}
\begin{enumerate}[label=\Alph*.]
    \item 2 times more
    \item 10 times more
    \item 20 times more
    \item 100 times more
    \item 1000 times more
\end{enumerate}
\textbf{Answer:} D \\
\textbf{Explanation:} Each pH unit represents a 10-fold change. A difference of 2 pH units means $10 \times 10 = 100$ times more $H^+$.

% ---------------------------------------
% 3.3 Buffers and Ocean Acidification (70-75)
% ---------------------------------------

\item \textbf{A substance that minimizes changes in the concentration of $H^+$ and $OH^-$ in a solution is called a:}
\begin{enumerate}[label=\Alph*.]
    \item Strong acid
    \item Strong base
    \item Buffer
    \item Solute
    \item Catalyst
\end{enumerate}
\textbf{Answer:} C \\
\textbf{Explanation:} Buffers resist changes in pH by accepting $H^+$ when it is in excess and donating $H^+$ when it is depleted.

\item \textbf{Most buffers in biological systems consist of:}
\begin{enumerate}[label=\Alph*.]
    \item A strong acid and a strong base.
    \item An acid-base pair that reversibly combines with hydrogen ions.
    \item Pure water and salt.
    \item Hydrophobic lipids.
    \item Carbon dioxide and oxygen.
\end{enumerate}
\textbf{Answer:} B \\
\textbf{Explanation:} Weak acids and their corresponding weak bases react reversibly to absorb or release $H^+$ to stabilize pH.

\item \textbf{Which buffer system is critical for maintaining the pH of human blood?}
\begin{enumerate}[label=\Alph*.]
    \item Acetic acid / acetate
    \item Carbonic acid / bicarbonate
    \item Hydrochloric acid / chloride
    \item Sulfuric acid / sulfate
    \item Ammonia / ammonium
\end{enumerate}
\textbf{Answer:} B \\
\textbf{Explanation:} Carbonic acid ($H_2CO_3$) dissociates to yield bicarbonate ($HCO_3^-$) and $H^+$, providing a crucial buffering system in the blood.

\item \textbf{Ocean acidification is primarily caused by oceans absorbing massive amounts of which gas?}
\begin{enumerate}[label=\Alph*.]
    \item Oxygen ($O_2$)
    \item Methane ($CH_4$)
    \item Carbon dioxide ($CO_2$)
    \item Nitrogen ($N_2$)
    \item Sulfur dioxide ($SO_2$)
\end{enumerate}
\textbf{Answer:} C \\
\textbf{Explanation:} About 25\% of human-generated $CO_2$ is absorbed by oceans, causing acidification.

\item \textbf{When $CO_2$ dissolves in seawater, it reacts with water to form:}
\begin{enumerate}[label=\Alph*.]
    \item Calcium carbonate
    \item Carbonic acid ($H_2CO_3$)
    \item Bicarbonate ($HCO_3^-$) exclusively
    \item Pure carbon
    \item Hydrochloric acid
\end{enumerate}
\textbf{Answer:} B \\
\textbf{Explanation:} Dissolved $CO_2$ reacts with $H_2O$ to form carbonic acid, which lowers the pH of the ocean.

\item \textbf{Ocean acidification threatens marine life such as corals and mollusks because it depletes the availability of:}
\begin{enumerate}[label=\Alph*.]
    \item Oxygen
    \item Hydronium ions
    \item Carbonate ions ($CO_3^{2-}$)
    \item Sodium ions
    \item Iron
\end{enumerate}
\textbf{Answer:} C \\
\textbf{Explanation:} The extra $H^+$ from carbonic acid combines with carbonate to form bicarbonate, reducing the carbonate available for marine organisms to build their calcium carbonate shells/skeletons.

\end{enumerate}

\end{document}
